\documentclass[12pt]{article}
\usepackage[utf8]{inputenc}
\usepackage[T1]{fontenc}
\usepackage{hyperref}
\title{Sardinian, Gallurese Language Resources for K-12 Curriculum Development}
\author{Compiled by Claude Sonnet 5.5 (claude-sonnet-5-5) for Andr\'as Kornai}
\date{}
\begin{document}
\maketitle

\section{Introduction}
Gallurese (ISO 639-3 \texttt{sdn}; Glottocode gall1276) is spoken in the Gallura region of northern Sardinia, Italy. It is endangered, and several sources treat it as distinct from central Sardinian varieties. Searches found no school textbook or dedicated dictionary for Gallurese specifically, but there are recent policy and orthography developments.

\section{Primary Educational Resources}
No school textbook adopted in schools was found. Adult and university courses exist:
\begin{itemize}
\item Press release on a 30-hour Gallurese course in Santa Teresa di Gallura, organised by Istituto Chircas. \url{https://santateresagallura-api.municipiumapp.it/s3/6523/allegati/allegati-luoghi/comunicato-stampa-santa-teresa-4173.pdf}
\item Poster of the same course. \url{https://santateresagallura-api.municipiumapp.it/s3/6523/allegati/allegati-luoghi/locandina-corso-di-gallurese-santa-teresa.pdf}
\end{itemize}

\section{Linguistic and Reference Materials}
\begin{itemize}
\item Wikipedia, Gallurese, which cites Rubattu's \emph{Elementi di grammatica gallurese}. \url{https://en.wikipedia.org/wiki/Gallurese}
\item Loporcaro (2008), study of the Gallurese of Luras. \url{https://ricerca.sns.it/retrieve/558c1fc9-8e2e-4611-a205-25dce3dd660b/122_Loporcaro_2008_Luras_post_print.pdf}
\item Loporcaro (2006), Luras, change and contact. \url{https://ricerca.sns.it/bitstream/11384/123843/1/108_Loporcaro_2006_Luras%20mutamento%20e%20contatto_open.pdf}
\item Omniglot, numbers in Gallurese. \url{https://omniglot.com/language/numbers/sardinian_gallurese.htm}
\item Wiktionary entry citing Rubattu's \emph{Dizionario universale della lingua di Sardegna} for Gallurese. \url{https://en.wiktionary.org/wiki/par%C3%A0ula}
\end{itemize}

\section{Religious and Community Resources}
No Gallurese scripture was found in the searches.

\section{Digital and Audio Resources}
None specific to Gallurese were found beyond the web references above.

\section{Curriculum Development Context}
The Sardinian regional government approved in 2024 an orthographic standard for Sassarese, Tabarchino and Gallurese. Regional law 22/2018 provides for teaching of the Sardinian languages in schools of every level.
\begin{itemize}
\item Regional council resolution relating to the orthographic standardisation. \url{https://delibere.regione.sardegna.it/protected/70103/0/def/ref/DBR70050/}
\item Regione Sardegna notice. \url{https://files.regione.sardegna.it/squidex/api/assets/redazionaleras/7e82d2b0-594e-48ed-b9e1-0682e2f41877/avviso.pdf}
\end{itemize}
The University of Sassari offers a 6-credit Gallurese laboratory, per a CV found in search. Teacher training and graded materials remain to be created.

\end{document}